\documentclass[11pt,a4paper]{article}
\usepackage{microtype}
\usepackage[margin=2.6cm]{geometry}
\usepackage{amsmath,amsthm,thmtools,thm-restate}
\usepackage{fontspec}
\usepackage{unicode-math}
\directlua{luaotfload.add_fallback("cfb", {"NotoSansMath:mode=harf;", "DejaVuSans:mode=harf;"})}
\setmainfont{Libertinus Serif}[RawFeature={fallback=cfb}]
\setmathfont{Latin Modern Math}
\setmonofont{Latin Modern Mono}[Scale=0.85,RawFeature={fallback=cfb}]
\AtBeginDocument{\let\setminus\smallsetminus}
% unicode-math defines \triangle as a protected \mathord{...}; a bare subscript _\triangle (as in the referee report) then fails. Make it a braced group.
\AtBeginDocument{\renewcommand{\triangle}{{\bigtriangleup}}}
\usepackage{enumitem}
\usepackage{longtable,booktabs,array,calc}
\usepackage{tcolorbox}
\usepackage{fancyhdr}
\usepackage[colorlinks=true,linkcolor=blue!50!black,citecolor=blue!50!black,urlcolor=blue!50!black]{hyperref}
\usepackage[capitalize,nameinlink]{cleveref}

\theoremstyle{plain}
\newtheorem{theorem}{Theorem}[section]
\newtheorem{lemma}[theorem]{Lemma}
\newtheorem{corollary}[theorem]{Corollary}
\newtheorem{proposition}[theorem]{Proposition}
\newtheorem{observation}[theorem]{Observation}
\theoremstyle{definition}
\newtheorem{definition}[theorem]{Definition}
\newtheorem{question}[theorem]{Question}
\newtheorem{conjecture}[theorem]{Conjecture}
\newtheorem{problem}[theorem]{Problem}
\theoremstyle{remark}
\newtheorem{remark}[theorem]{Remark}
\newtheorem*{proofidea}{Idea of proof}
\newcommand{\fullproofref}[1]{\,\hyperref[#1]{\textit{[full proof]}}}
\providecommand{\tightlist}{\setlength{\itemsep}{0pt}\setlength{\parskip}{0pt}}

\newcommand{\Tstar}{T_{\ast}}
\newcommand{\TT}{\mathrm{TT}}
\newcommand{\fin}{\subseteq^{\ast}}
\newcommand{\symdiff}{\mathbin{\triangle}}
\newcommand{\NN}{\mathbb N}
\newcommand{\QQ}{\mathbb Q}

\pagestyle{fancy}\fancyhf{}
\fancyhead[L]{\small\itshape Candidate result 2506.08810\_\_03 --- machine-generated proof, awaiting human review}
\fancyhead[R]{\small\thepage}
\renewcommand{\headrulewidth}{0.2pt}

\title{The five-vertex tournament $C_3[\TT_2,\TT_2,1]$ is a counterexample\\ to Conjecture~24 of Bonamy, Groenland, Johnston, Morrison and Scott}
\author{\href{https://github.com/graph-theory-AI}{github.com/graph-theory-AI}\\[2pt] \normalsize Machine-generated note}
\date{17 September 2026}

\begin{document}
\maketitle

\begin{abstract}
Bonamy, Groenland, Johnston, Morrison and Scott proved that for every finite graph $H$ other than a clique or an independent set there is a countably infinite $H$-free graph every locally finite perturbation of which contains an induced copy of $H$, and they conjectured the tournament analogue (Conjecture~24 of arXiv:2506.08810): for every finite non-transitive tournament $T$ there is a countably infinite $T$-free tournament $S_T$ every locally finite perturbation of which contains $T$. We show that the five-vertex tournament $\Tstar=C_3[\TT_2,\TT_2,1]$, a directed triangle two of whose vertices are replaced by transitive tournaments on two vertices, is a counterexample. The reason is structural: every countable $\Tstar$-free tournament is obtained from a transitive tournament by reversing a locally finite set of arcs, and transitive tournaments are themselves $\Tstar$-free; hence every countably infinite $\Tstar$-free tournament has a nonempty locally finite perturbation that is still $\Tstar$-free. The single-arc variant, Problem~23 of the same paper, is not addressed.

\medskip\noindent\small
This note was produced by AI within the \href{https://github.com/graph-theory-AI/Graph-Theory-LLM-Proofs}{Graph Theory AI} project:
the argument was found by a language model and audited by an adversarial LLM
referee, and it has not been checked by a human mathematician.
\Cref{sec:provenance} records the full provenance.
\end{abstract}

\section{Introduction}

\paragraph{Tournaments.}
A \emph{tournament} is an orientation of a complete graph; we write $x\to y$ when the arc between $x$ and $y$ is directed from $x$ to $y$, and $O(x)=N^+(x)$ for the set of out-neighbours of $x$. For sets $X,Y$ of vertices, $X\to Y$ means $x\to y$ for all $x\in X$, $y\in Y$. A tournament is \emph{transitive} if it contains no directed triangle; equivalently its arcs are $x\to y$ for $x<y$ in some linear order of its vertex set. $\TT_n$ denotes the transitive tournament on $n$ vertices and $C_3$ the directed triangle. For tournaments $T_0,T_1,T_2$, the \emph{composition} $C_3[T_0,T_1,T_2]$ is the tournament on the disjoint union $V(T_0)\cup V(T_1)\cup V(T_2)$ that induces $T_i$ on $V(T_i)$ and has $V(T_0)\to V(T_1)\to V(T_2)\to V(T_0)$; a single vertex is written $1$. A tournament $S$ is \emph{$T$-free} if no set of $|V(T)|$ vertices of $S$ induces a tournament isomorphic to $T$. Every subtournament of a tournament is induced, so ``induced copy of $T$'' and ``copy of $T$'' coincide for tournaments. A tournament is \emph{countable} if its vertex set is finite or countably infinite; the conjecture below concerns countably infinite ones.

\paragraph{Perturbations.}
For tournaments $S,S'$ on the same vertex set, the \emph{change graph} $D(S,S')$ is the undirected graph on $V(S)$ whose edges are the pairs whose orientation differs between $S$ and $S'$. Bonamy, Groenland, Johnston, Morrison and Scott \cite[Section~3]{BGJMS25} define the operative notion for graphs as follows:
\begin{quote}\small
For a pair of distinct vertices $u,v\in V(G)$, perturbing the pair $uv$ is the operation which removes $uv$ from $E(G)$ if it is present and adds it if it is not present. A locally finite perturbation of $G$ is a graph obtained from $G$ by perturbing arbitrarily many pairs (at least one, possibly infinitely many) under the constraint that for any vertex $v$, the number of perturbed pairs involving $v$ is finite.
\end{quote}
Their Introduction phrases the same notion as ``perturbing a nonempty set of pairs such that every vertex is in a finite number of perturbed pairs; equivalently, we can consider graphs $G'$ such that $E(G)\symdiff E(G')$ is the edge set of a locally finite graph'' \cite[Section~1]{BGJMS25}, which is the change-graph formulation used below. In their concluding section they add that ``for tournaments the natural operation is to reverse the direction of edges'' \cite[Section~8]{BGJMS25}. Accordingly, a \emph{locally finite perturbation} of a tournament $S$ is a tournament $S'$ on $V(S)$ such that $D(S,S')$ is nonempty and locally finite (every vertex lies in finitely many changed pairs). We stress that nonemptiness is part of the paper's definition (``at least one''), not a convention introduced here; the paper itself notes that ``the locally finite perturbations enforce that $S_T$ is infinite''.

\paragraph{The problem.}
The main result of \cite{BGJMS25} (their Theorem~1) is that for every finite graph $H$ which is not a clique or an independent set there is a countably infinite $H$-free graph $G_H$ such that every locally finite perturbation of $G_H$ has an induced copy of $H$. Section~8 of \cite{BGJMS25} asks for analogues in other structures and poses, for tournaments, first the single-arc problem and then the conjecture we refute. Both are quoted verbatim.

\begin{problem}[{\cite[Problem~23]{BGJMS25}}]\label{prob:23}
For which finite tournaments $H$, does there exist a countable tournament $G$ that does not contain $H$ as subtournament, yet any tournament obtained from $G$ by changing the direction of a single arc does contain $H$?
\end{problem}

\begin{conjecture}[{\cite[Conjecture~24]{BGJMS25}}]\label{conj:24}
Let $T$ be a finite tournament. If $T$ is not transitive, then there exists a countably infinite $T$-free tournament $S_{T}$ such that every locally finite perturbation of $S_{T}$ has an induced copy of $T$.
\end{conjecture}

The conjecture is introduced with the sentence ``Perhaps the following analogue of Theorem~1 is true for the `locally finite perturbation' variant of this question.'' Its only hypothesis is that $T$ is not transitive, and the paper explains why this hypothesis is necessary: every tournament on $2^n$ vertices contains a transitive subtournament on $n$ vertices, so a transitive $T$ cannot be avoided by an infinite tournament at all. There is no restriction on the order, connectivity or modular structure of $T$. The catalog entry for this problem (id \texttt{2506.08810\_\_03}) and the referee's own searches found no published work addressing \cref{conj:24}; the referee also found that the journal version (Canadian Journal of Mathematics, published online 24 March 2026) contains no retraction of, or note on, the tournament conjecture.

\paragraph{The counterexample.}
Let $\Tstar=C_3[\TT_2,\TT_2,1]$ be the tournament on the five vertices $a_0,a_1,b_0,b_1,c$ with
\[
 a_0\to a_1,\qquad b_0\to b_1,\qquad \{a_0,a_1\}\to\{b_0,b_1\}\to\{c\}\to\{a_0,a_1\},
\]
that is, with the ten arcs
\begin{gather*}
 a_0\to a_1,\quad a_0\to b_0,\quad a_0\to b_1,\quad a_1\to b_0,\quad a_1\to b_1,\\
 b_0\to b_1,\quad b_0\to c,\quad b_1\to c,\quad c\to a_0,\quad c\to a_1 .
\end{gather*}
The out-degrees of $(a_0,a_1,b_0,b_1,c)$ are $(3,2,2,1,2)$. $\Tstar$ has exactly four directed triangles, namely $c\to a_i\to b_j\to c$ for $i,j\in\{0,1\}$, all through $c$; in particular it is not transitive, and the two triangles $c\to a_0\to b_0\to c$ and $c\to a_1\to b_1\to c$ share no arc. Both $\{a_0,a_1\}$ and $\{b_0,b_1\}$ are nontrivial modules, so $\Tstar$ is not prime; $\Tstar$ is strongly connected but $\Tstar-c$ is not; and the minimum number of arcs whose reversal makes $\Tstar$ transitive is $2$.

\begin{restatable}[Main theorem]{theorem}{thmmain}\label{thm:main}
Let $S$ be a countably infinite $\Tstar$-free tournament. Then there is a $\Tstar$-free tournament $S'\ne S$ on $V(S)$ such that the change graph $D(S,S')$ is locally finite. Equivalently: every countably infinite $\Tstar$-free tournament has a locally finite perturbation (in the sense of \cite{BGJMS25}) that is still $\Tstar$-free.
\end{restatable}

\begin{corollary}\label{cor:main}
\Cref{conj:24} is false. The tournament $T=\Tstar$ is finite and not transitive, yet no countably infinite $\Tstar$-free tournament $S_{T}$ has the property that every locally finite perturbation of $S_T$ contains an induced copy of $\Tstar$.
\end{corollary}
\begin{proof}
$\Tstar$ is a finite tournament containing a directed triangle, so it satisfies the only hypothesis of \cref{conj:24}. Let $S_T$ be any countably infinite $\Tstar$-free tournament. \Cref{thm:main} supplies $S'\ne S_T$ on the same vertex set with $D(S_T,S')$ locally finite; since $S'\ne S_T$ the set of reversed pairs is nonempty, so $S'$ is a locally finite perturbation of $S_T$ in the paper's sense, and $S'$ is $\Tstar$-free, i.e.\ it has no induced copy of $\Tstar$. Hence $S_T$ does not have the conjectured property.
\end{proof}

\begin{proofidea}
Two finite facts about $\Tstar$ drive the argument, and the referee checked both by exhaustive enumeration. First, if an arc $x\to y$ of a $\Tstar$-free tournament admits a vertex $p$ with $x\to p\to y$, then at most two vertices $z$ satisfy $y\to z\to x$ (\cref{lem:Z}): any three such $z$ together with $x,y,p$ contain a copy of $\Tstar$. So an arc lying in infinitely many directed triangles has no such $p$, and arcs with no such $p$ form an undirected graph of maximum degree at most $2$ (\cref{lem:Pempty}). Hence the graph $J(S)$ of arcs in infinitely many directed triangles has maximum degree at most $2$, and a general criterion (\cref{lem:criterion}) shows that a countable tournament with locally finite $J(S)$ differs from some transitive tournament $L$ by a locally finite set of reversals (\cref{thm:structure}). If $S$ is not transitive then $L\ne S$, and $L$ is $\Tstar$-free because $\Tstar$ is not transitive. Second, if $S$ is transitive, reversing a single arc creates only directed triangles through that arc, whereas $\Tstar$ contains two arc-disjoint directed triangles; so the result is still $\Tstar$-free (\cref{lem:onearc}).\fullproofref{pf:main}
\end{proofidea}

\section{Structure of \texorpdfstring{$\Tstar$}{T*}-free tournaments}\label{sec:structure}

For an arc $x\to y$ of a tournament $S$ put
\[
 Z(x,y)=\{z: y\to z\to x\},\qquad P(x,y)=\{p: x\to p\to y\}.
\]
$Z(x,y)$ is the set of vertices completing $x\to y$ to a directed triangle, and $P(x,y)$ the set of vertices that would do so if the arc $xy$ were reversed. Note that $O(y)\setminus O(x)=Z(x,y)$: indeed $x\notin O(y)$ because $x\to y$, so a vertex $z\in O(y)\setminus O(x)$ is distinct from $x$ and $y$, satisfies $y\to z$ and, not being an out-neighbour of $x$, satisfies $z\to x$. Let $J(S)$ be the undirected graph on $V(S)$ in which $xy$ is an edge precisely when the arc between $x$ and $y$ lies in infinitely many directed triangles, i.e.\ when $Z(x,y)$ (for $x\to y$) is infinite.

\begin{restatable}[Criterion for locally finite distance from a transitive tournament]{lemma}{lemcriterion}\label{lem:criterion}
Let $S$ be a countable tournament. There is a transitive tournament $L$ on $V(S)$ with $D(S,L)$ locally finite if and only if $J(S)$ is locally finite.
\end{restatable}
\begin{proofidea}
Necessity: if $xy\notin E(D(S,L))$ then every directed triangle $xyz$ of $S$ must have $xz$ or $yz$ in $D(S,L)$, else it survives into $L$; so $J(S)\subseteq D(S,L)$. Sufficiency: call $x\sim y$ if $O(x)\symdiff O(y)$ is finite. On the classes, ``$O(y)\setminus O(x)$ finite'' defines a partial order; take a linear extension (Szpilrajn \cite{Sz30}; for countably many classes an explicit insertion argument suffices). By $O(y)\setminus O(x)=Z(x,y)$, any arc between different classes that points against this order lies in $J(S)$. Inside a class $C$, fix $r\in C$ and split $C$ into the in-neighbours $C^-$ and out-neighbours $C^+$ of $r$; finiteness of $O(x)\symdiff O(r)$ gives finite out-degrees inside $C^-$, finite in-degrees inside $C^+$, and only locally finitely many arcs from $C^+$ to $C^-$. Ordering $C^-$ in reverse enumeration order, then $r$, then $C^+$ in enumeration order, makes the discrepancies inside $C$ locally finite. Only the sufficiency direction is used below.\fullproofref{pf:criterion}
\end{proofidea}

\begin{restatable}{lemma}{lemZ}\label{lem:Z}
Let $S$ be $\Tstar$-free and let $x\to y$ be an arc with $P(x,y)\ne\varnothing$. Then $|Z(x,y)|\le 2$.
\end{restatable}
\begin{proofidea}
Take $p\in P(x,y)$ and suppose $z_1,z_2,z_3\in Z(x,y)$. All five vertices are distinct, and two of the $z_i$, say $z_1,z_2$, are oriented the same way relative to $p$. If $p\to z_1,z_2$ then $\{p,y\}\to\{z_1,z_2\}\to\{x\}\to\{p,y\}$ with $p\to y$; if $z_1,z_2\to p$ then $\{z_1,z_2\}\to\{x,p\}\to\{y\}\to\{z_1,z_2\}$ with $x\to p$. Either way the five vertices induce $C_3[\TT_2,\TT_2,1]=\Tstar$ (a two-vertex block is $\TT_2$ whichever way its arc points).\fullproofref{pf:Z}
\end{proofidea}

\begin{restatable}{lemma}{lemPempty}\label{lem:Pempty}
In any tournament, the arcs $x\to y$ with $P(x,y)=\varnothing$ form an undirected graph of maximum degree at most $2$: every vertex has at most one such out-arc and at most one such in-arc.
\end{restatable}
\begin{proofidea}
If $x\to y$ and $x\to z$ both had empty $P$, then $y\to z$ would put $y\in P(x,z)$ and $z\to y$ would put $z\in P(x,y)$. In-arcs are symmetric.\fullproofref{pf:Pempty}
\end{proofidea}

\begin{restatable}[Structure theorem]{theorem}{thmstructure}\label{thm:structure}
Let $S$ be a $\Tstar$-free tournament. Then $J(S)$ has maximum degree at most $2$. Consequently, if $S$ is countable, there is a transitive tournament $L$ on $V(S)$ such that $D(S,L)$ is locally finite: every countable $\Tstar$-free tournament is obtained from a transitive tournament by reversing a locally finite set of arcs.
\end{restatable}
\begin{proofidea}
An edge $xy$ of $J(S)$ with $x\to y$ has $Z(x,y)$ infinite, hence of size at least $3$, so $P(x,y)=\varnothing$ by \cref{lem:Z}. Thus $J(S)$ is a subgraph of the graph of \cref{lem:Pempty}, and \cref{lem:criterion} applies.\fullproofref{pf:structure}
\end{proofidea}

\begin{restatable}{lemma}{lemonearc}\label{lem:onearc}
Let $L$ be a transitive tournament and let $L'$ be obtained from $L$ by reversing a single arc. Then $L'$ is $\Tstar$-free.
\end{restatable}
\begin{proofidea}
Every directed triangle of $L'$ contains the reversed arc, since otherwise it would be a directed triangle of $L$. But $\Tstar$ contains two arc-disjoint directed triangles, so no copy of $\Tstar$ fits into $L'$.\fullproofref{pf:onearc}
\end{proofidea}

The structure theorem is one-directional: a locally finite perturbation of a transitive tournament may well contain $\Tstar$. Indeed $\Tstar$ itself is obtained from $\TT_5$ with the order $a_0<a_1<b_0<b_1<c$ by reversing the two arcs $a_0c$ and $a_1c$, and no single reversal suffices (\cref{lem:onearc}), so the inversion number of $\Tstar$ is exactly $2$. What \cref{thm:structure} says is that nothing outside the locally finite perturbations of transitive tournaments is $\Tstar$-free.

\section{Remarks}

\begin{remark}[Interpretation audit]\label{rem:interp}
Because the object is tiny, the substance of the referee's audit was to check that \cref{conj:24} as posed really admits it. The referee fetched Section~8 of arXiv:2506.08810v3 and found the conjecture's statement to agree character for character with the one attacked (\cref{conj:24} above); it fetched the definition of ``locally finite perturbation'' quoted in the Introduction and found the writeup's reading, a nonempty set of reversed arcs with every vertex in finitely many of them, to be the paper's own. Three consequences were recorded. (i)~The nonemptiness caveat the writeup states is not an assumption it imposed but the paper's convention (``at least one''); had the writeup instead used the empty perturbation, the disproof would be a strawman, and it does not. This is why the case of a transitive $S$ in \cref{thm:main} needs the separate \cref{lem:onearc}: the transitive tournament $L$ produced by \cref{thm:structure} coincides with $S$ there, so a genuine reversal must be exhibited. (ii)~``Induced copy'' means ``subtournament'', so there is no induced-versus-weak ambiguity to exploit. (iii)~The conjecture has no order threshold and no primeness or connectivity hypothesis; its only hypothesis, non-transitivity, is satisfied by $\Tstar$, which has four directed triangles. The referee concluded that a reasonable author of \cref{conj:24} would accept the argument as a refutation of the conjecture as posed.
\end{remark}

\begin{remark}[Why the conjecture admits a five-vertex counterexample]\label{rem:why}
The referee was instructed to explain why five authors would have missed a five-vertex counterexample, and gave three answers, which we reproduce as its assessment. (i)~\Cref{conj:24} is a speculation in the open-problems section of a paper about graphs, introduced with ``Perhaps the following analogue of Theorem~1 is true''; nothing in the paper indicates a small-tournament check, and the paper's ancillary code (\texttt{small\_graphs.py}, \texttt{theorem\_10.py}) concerns the graph results. (ii)~The intuition behind Theorem~1 of \cite{BGJMS25} is that once the degenerate targets (cliques and independent sets) are excluded, the space of $H$-free graphs has ``points that are isolated in a very strong sense''; the natural transfer replaces the degenerate targets by transitive tournaments. For $\Tstar$ this fails for a tournament-specific reason: the obstruction is not that $\Tstar$ is hard to hit but that it is so easy to avoid that, by \cref{thm:structure}, the whole space of countable $\Tstar$-free tournaments collapses onto the locally finite perturbations of transitive tournaments; and transitive tournaments are themselves $\Tstar$-free, so this space has no isolated point at all. (iii)~The machinery of \cite{BGJMS25} replaces vertices by infinite cliques or independent sets; the tournament analogue replaces vertices by infinite transitive blocks. But $\Tstar=C_3[\TT_2,\TT_2,1]$ is exactly a directed triangle with two vertices blown up into transitive blocks of size $2$, so any blow-up of a directed triangle with two blocks of size at least $2$ already contains $\Tstar$, and the technique cannot get started for this target. The referee also measured how special $\Tstar$ is: among all tournaments on $4$, $5$ and $6$ vertices up to isomorphism ($4$, $12$ and $56$ of them), $\Tstar$ is the \emph{only} one for which both finite engines of the proof work, namely all $64$ completions of the configuration in \cref{lem:Z} contain $T$, and $T$ has two arc-disjoint directed triangles. The counterexample is thus a narrow object rather than one that should have been obvious.
\end{remark}

\begin{remark}[Scope: Problem~23 is untouched]\label{rem:p23}
\Cref{thm:main} says nothing about \cref{prob:23} for $\Tstar$. In the non-transitive case its perturbation $L$ in general reverses infinitely many arcs, and in the transitive case \cref{lem:onearc} shows only that a transitive tournament is not a solution to \cref{prob:23} for $H=\Tstar$. Whether some countable $\Tstar$-free tournament becomes non-$\Tstar$-free under every single arc reversal remains open, as the writeup itself states. Likewise the writeup proposes no repaired conjecture; the question for which non-transitive $T$ the conclusion of \cref{conj:24} holds is open. The earlier attempt discussed in \cref{rem:prior} proves it for several classes.
\end{remark}

\begin{remark}[Relation to the earlier attempt]\label{rem:prior}
This target was a second attempt. The first, by \texttt{gpt-5.6-sol} (\texttt{attacks/2506.08810\_\_03/output.md}, verdict ``partial'', unverified), proved the conclusion of \cref{conj:24} for every tournament one arc reversal away from transitive, for the cyclic tournament $C_5$ via a $C_3$-lexicographic construction, and for tournaments that stay strongly connected after deleting any two vertices and have no proper transitive module (including $C_7$), and it named low-connectivity tournaments with nontrivial modular structure as the untreated region. The present argument reuses nothing from it. The referee checked that the two are consistent: $\Tstar$ needs $2$ reversals to become transitive (exhaustive over all $120$ orderings), has the nontrivial modules $\{a_0,a_1\}$ and $\{b_0,b_1\}$, and $\Tstar-c$ is not strongly connected; so none of the earlier positive results covers $\Tstar$, which falls exactly into the region that attempt flagged as open. Had the inversion number of $\Tstar$ been $1$, one of the two documents would be wrong.
\end{remark}

\begin{remark}[Computational checks by the referee]\label{rem:comp}
The referee's scripts (pure Python, in the repository folder \nolinkurl{verification_astra/scripts/2506.08810__03/}) established the following; none failed.
\begin{itemize}\tightlist
\item \emph{The object.} $\Tstar$ is a tournament with out-degree sequence $(3,2,2,1,2)$ on $(a_0,a_1,b_0,b_1,c)$,\footnote{The referee's report lists this as $(3,2,2,2,1)$ in the same order; $b_1$ has out-degree $1$ (it beats only $c$) and $c$ has out-degree $2$, so the report's two last entries are transposed. The multiset, and everything built on it, is unaffected. We re-derived all invariants in this paragraph.} four directed triangles all through $c$, a pair of arc-disjoint directed triangles, empty intersection of the arc sets of all four triangles, inversion number $2$, nontrivial modules $\{a_0,a_1\}$ and $\{b_0,b_1\}$, strongly connected with $\Tstar-c$ not strongly connected.
\item \emph{\Cref{lem:Z} is a complete enumeration.} The configuration $x\to y$, $p\in P(x,y)$, $z_1,z_2\in Z(x,y)$ with $z_1,z_2$ alike relative to $p$ has $4$ completions, all isomorphic to $\Tstar$; the six-vertex configuration with $z_1,z_2,z_3\in Z(x,y)$ has $2^3\cdot 2^3=64$ completions and $0$ of them are $\Tstar$-free.
\item \emph{\Cref{lem:Pempty}.} Three-vertex configurations with two $P$-empty out-arcs at a vertex: $0$; with two $P$-empty in-arcs: $0$.
\item \emph{Finite shadow of \cref{thm:structure}.} Over all labelled tournaments on $n$ vertices, the arcs lying in at least three directed triangles of a $\Tstar$-free tournament always form a graph of maximum degree at most $2$ with $P=\varnothing$ throughout:
\begin{center}\small
\begin{tabular}{@{}rrrrl@{}}\toprule
$n$ & labelled tournaments & $\Tstar$-free & max.\ degree of the $\ge3$-triangle graph & such an arc with $P\ne\varnothing$\\\midrule
4 & 64 & 64 & 0 & none\\
5 & 1\,024 & 904 & 1 & none\\
6 & 32\,768 & 17\,648 & 1 & none\\
7 & 2\,097\,152 & 434\,352 & 2 & none\\\bottomrule
\end{tabular}
\end{center}
The bound $2$ is attained at $n=7$, so the constant in \cref{lem:Pempty} is tight. A further $24$ $\Tstar$-free tournaments on $8$, $9$, $10$ and $12$ vertices found by randomised local search satisfy both conclusions.
\item \emph{\Cref{lem:onearc}.} Reversing each of the $\binom n2$ arcs of $\TT_n$ for $n=5,\dots,9$ ($110$ cases) creates $0$ copies of $\Tstar$.
\item \emph{Infinite examples.} On finite windows of eight explicit countable tournaments, the four that contain $\Tstar$ (among them $C_3[\TT,\TT,\TT]$, the circular tournament on $\mathbb Z_{17}$, and the earlier attempt's $C_3$-lexicographic power) are as \cref{thm:structure} requires, and the four $\Tstar$-free ones ($\TT_\omega[C_3]$; $\TT_\omega$ with the pairs $(2t,2t+1)$ reversed; $\TT_\omega$ with all pairs $(i,i+1)$ reversed; evens dominating odds minus a matching) have triangle graph of maximum degree at most $1$ on the window. For these four the recipe of \cref{lem:criterion} was worked out by hand and its change graph measured on windows of $20$, $60$, $120$ and $240$ vertices: the maximum degree stays at $1$, $1$, $2$ and $1$ respectively, constant in the window size. The referee notes that ``$O(x)\symdiff O(y)$ finite'' has no faithful finite-window proxy, which is why this check was done by hand rather than automatically.
\item \emph{Uniqueness on at most six vertices}, as described in \cref{rem:why}.
\end{itemize}
\end{remark}

\begin{remark}[Use of choice]\label{rem:choice}
The sufficiency direction of \cref{lem:criterion} orders the $\sim$-classes by a linear extension of a partial order and enumerates each class. In general the existence of a linear extension is Szpilrajn's theorem \cite{Sz30}; for a countable tournament there are countably many classes and the extension can be built by inserting the classes one at a time (\cref{app:proofs}), so the argument uses no choice beyond fixing enumerations of a countable set. The writeup did not name Szpilrajn; the referee called this cosmetic.
\end{remark}

\section{Provenance}\label{sec:provenance}

The mathematics in this note was produced without human intervention, in three stages.
(1)~The counterexample and its proof were found by the language model \texttt{gpt-6-astra} (reasoning effort max, flex tier) in a single autonomous pass started on 2026-09-17, in the retry leg of the campaign (catalog id \texttt{2506.08810\_\_03}, leg \texttt{attacks\_retry}; writeup in \texttt{attacks\_retry/2506.08810\_\_03/output.md}). This was a second attempt at the problem: the prompt included the earlier attempt by \texttt{gpt-5.6-sol} (\texttt{attacks/2506.08810\_\_03/output.md}, verdict ``partial'', no referee report), which is summarised in \cref{rem:prior}. The new writeup states that it uses no result of the earlier one, and the referee confirmed that nothing is inherited; the partial results of \cref{rem:prior} are the earlier model's and are not claimed here.
(2)~An independent adversarial referee agent (\texttt{claude-opus-5}) reviewed the writeup on 2026-09-17. It fetched arXiv:2506.08810v3 and checked the statement of Conjecture~24 and the definition of ``locally finite perturbation'' verbatim against the writeup's reading, re-derived every step (in particular the three ``it follows that'' bullets in the sufficiency direction of \cref{lem:criterion} and the implicit step that a cross-class arc contradicting the linear extension lies in $J(S)$), verified \cref{lem:Z}, \cref{lem:Pempty} and \cref{lem:onearc} by complete enumeration, checked the degree bound over all $2\,097\,152$ labelled tournaments on $7$ vertices, checked consistency with the earlier attempt, and searched for prior or concurrent work, finding none. Its verdict was \textsc{confirmed}, high confidence. Its report is reproduced verbatim in \cref{app:referee}.
(3)~On 2026-09-17 the writeup was rewritten into the present note by \texttt{claude-fable-5-1}; the compilation was repaired and the quotations from the paper re-checked on 2026-09-18. The deviations from the writeup are the following, and no other mathematical content was changed.
\begin{itemize}\tightlist
\item The text was reorganised with full proofs in \cref{app:proofs}. The writeup's Corollary~4 is stated as \cref{thm:structure}, and its ``Case~2'' is isolated as \cref{lem:onearc}.
\item The explicit arc list of $\Tstar$ and its invariants (out-degrees, directed triangles, modules, connectivity, inversion number) were added; they come from the referee's step~1 and were recomputed for this note, correcting the labelling slip noted in \cref{rem:comp}.
\item In \cref{lem:criterion}: the existence of a linear extension is attributed to Szpilrajn \cite{Sz30} and an explicit construction for countably many classes is supplied (referee caveat~3); the step that an arc between classes pointing against the linear extension, and not merely against the partial order, lies in $J(S)$ is spelled out (referee step~S4); the three bullets on $C^-$ and $C^+$ are derived in full (referee step~S5); the well-definedness and transitivity of the order on classes are written out.
\item In \cref{lem:Z} the distinctness of the five vertices is made explicit (referee step~S7).
\item The verbatim quotations of the paper's definition of locally finite perturbation (from its Section~3 and, in the equivalent change-graph form, its Section~1), of Problem~23, of the sentence introducing Conjecture~24 and of the remark that the perturbations force $S_T$ to be infinite were added. The referee fetched Conjecture~24 and the definition from arXiv:2506.08810v3; all five passages were re-fetched from the same source and confirmed character for character on 2026-09-18. The referee report locates the definition in Section~2 of the paper; it is in Section~3 (``Notation and definitions''), and the note cites it there.
\item The paragraph following \cref{thm:structure}, observing that $\Tstar$ is $\TT_5$ with the two arcs $a_0c$ and $a_1c$ reversed and hence has inversion number exactly $2$, was added; the value $2$ is the referee's exhaustive computation, and the explicit pair of arcs was checked by hand against the arc list in the Introduction.
\item The bibliographic data of the journal version (DOI 10.4153/S0008414X26102132, published online 24 March 2026) were taken from the publisher's page on 2026-09-18.
\item \Cref{rem:interp,rem:why,rem:prior,rem:comp,rem:choice} are taken from the referee report; they contain no claim of the writeup's.
\end{itemize}
\textbf{No human mathematician has checked this argument.} We circulate it for human review, not as a claim of a settled result.

\appendix

\section{Full proofs}\label{app:proofs}

Throughout, for an arc $x\to y$ of the tournament under consideration, $Z(x,y)=\{z: y\to z\to x\}$, $P(x,y)=\{p: x\to p\to y\}$, $O(x)=N^+(x)$, and $X\fin Y$ means that $X\setminus Y$ is finite. We use repeatedly the identity
\begin{equation}\label{eq:Z}
 O(y)\setminus O(x)=Z(x,y)\qquad\text{whenever }x\to y,
\end{equation}
proved at the start of \cref{sec:structure}: $x\notin O(y)$ since $x\to y$, so $z\in O(y)\setminus O(x)$ means $z\ne x,y$, $y\to z$ and $z\to x$.

\phantomsection\label{pf:criterion}
\lemcriterion*
\begin{proof}
\emph{Necessity.} Let $L$ be transitive on $V(S)$ with $D=D(S,L)$ locally finite. We show $J(S)\subseteq D$, which gives local finiteness of $J(S)$. Let $xy\notin E(D)$, so $xy$ has the same orientation in $S$ and $L$, and let $z$ be such that $x,y,z$ span a directed triangle of $S$. If neither $xz$ nor $yz$ were in $E(D)$, all three arcs of this triangle would have the same orientation in $L$, and $L$ would contain a directed triangle, contradicting transitivity. Hence $z\in N_D(x)\cup N_D(y)$, a finite set. So $xy$ lies in only finitely many directed triangles of $S$, i.e.\ $xy\notin E(J(S))$.

\emph{Sufficiency.} Assume $J(S)$ is locally finite.

\emph{Step 1: classes.} Define $x\sim y$ iff $O(x)\symdiff O(y)$ is finite. This is reflexive and symmetric, and transitive because $O(x)\symdiff O(z)\subseteq (O(x)\symdiff O(y))\cup(O(y)\symdiff O(z))$. Let $\mathcal C$ be the set of $\sim$-classes; it is countable.

\emph{Step 2: a partial order on classes.} For distinct classes $C,D$ declare $C\prec D$ if $O(y)\fin O(x)$ for some (equivalently every) $x\in C$, $y\in D$. This is independent of the representatives: if $x'\sim x$ and $y'\sim y$ then $O(y')\setminus O(x')\subseteq (O(y')\setminus O(y))\cup(O(y)\setminus O(x))\cup(O(x)\setminus O(x'))$, and the first and third sets are finite. It is transitive, since $O(z)\setminus O(x)\subseteq (O(z)\setminus O(y))\cup(O(y)\setminus O(x))$. It is antisymmetric: if $C\prec D$ and $D\prec C$ then both $O(y)\setminus O(x)$ and $O(x)\setminus O(y)$ are finite, so $x\sim y$ and $C=D$. Thus $\prec$ is a strict partial order on $\mathcal C$.

\emph{Step 3: a linear extension.} Let $<$ be a linear order on $\mathcal C$ extending $\prec$. Such an order exists by Szpilrajn's theorem \cite{Sz30}; since $\mathcal C$ is countable we can also construct it directly. Enumerate $\mathcal C=\{C_1,C_2,\dots\}$ and build linear orders $<_n$ on $\{C_1,\dots,C_n\}$ extending $\prec$, each extending the previous one. Given $<_n$, let $\mathcal P=\{C_i: i\le n,\ C_i\prec C_{n+1}\}$ and $\mathcal Q=\{C_j: j\le n,\ C_{n+1}\prec C_j\}$. For $C_i\in\mathcal P$ and $C_j\in\mathcal Q$ transitivity gives $C_i\prec C_j$, hence $C_i<_n C_j$; so every element of $\mathcal P$ is $<_n$-below every element of $\mathcal Q$. Insert $C_{n+1}$ immediately after the $<_n$-largest element of $\mathcal P$ (at the beginning if $\mathcal P=\varnothing$). The result $<_{n+1}$ is a linear order placing $C_{n+1}$ above $\mathcal P$ and below $\mathcal Q$, hence extending $\prec$. The union $<$ of the $<_n$ is a linear order on $\mathcal C$ extending $\prec$.

\emph{Step 4: arcs between classes.} We claim: if $x\in C$, $y\in D$, $C\ne D$, $C<D$ and $y\to x$ in $S$, then $xy\in E(J(S))$. Suppose not: then $Z(y,x)$ is finite, so by \eqref{eq:Z} (applied to the arc $y\to x$) $O(x)\setminus O(y)$ is finite, i.e.\ $O(x)\fin O(y)$, which by definition means $D\prec C$, hence $D<C$, contradicting $C<D$. So every pair between distinct classes whose orientation in $S$ disagrees with $<$ is an edge of $J(S)$.

\emph{Step 5: inside a class.} Fix a class $C$ and $r\in C$, and put $C^-=\{x\in C: x\to r\}$, $C^+=\{x\in C: r\to x\}$, so $C=C^-\cup\{r\}\cup C^+$. For $x\in C$ the set $F(x)=O(x)\symdiff O(r)$ is finite. We verify three facts.
\begin{enumerate}[label=(\roman*)]\tightlist
\item If $x,y\in C^-$ and $x\to y$, then $y\in O(x)$ and $y\notin O(r)$ (as $y\to r$), so $y\in O(x)\setminus O(r)\subseteq F(x)$. Hence every $x\in C^-$ has finitely many out-neighbours in $C^-$.
\item If $x,y\in C^+$ and $y\to x$, then $y\in O(r)$ and $y\notin O(x)$, so $y\in O(r)\setminus O(x)\subseteq F(x)$. Hence every $x\in C^+$ has finitely many in-neighbours in $C^+$.
\item If $x\in C^-$, $y\in C^+$ and $y\to x$, then $y\in O(r)\setminus O(x)\subseteq F(x)$ and $x\in O(y)\setminus O(r)\subseteq F(y)$ (as $x\to r$, $x\notin O(r)$). Hence the bipartite graph of pairs $x\in C^-$, $y\in C^+$ with $y\to x$ has all degrees finite.
\end{enumerate}
Enumerate $C^-=\{u_1,u_2,\dots\}$ and $C^+=\{v_1,v_2,\dots\}$ (finite or infinite enumerations, without repetition) and order $C$ as
\begin{equation}\label{eq:classorder}
 \cdots<u_3<u_2<u_1<r<v_1<v_2<v_3<\cdots,
\end{equation}
i.e.\ $C^-$ in reverse enumeration order, then $r$, then $C^+$ in enumeration order. Let $L_C$ be the transitive tournament on $C$ given by this order ($a\to b$ iff $a<b$). We count the pairs inside $C$ on which $S$ and $L_C$ disagree at a fixed vertex. At $u_i$: with $r$ none, since $u_i\to r$ in both; with $u_j$ for $j<i$ at most $i-1$ pairs; with $u_j$ for $j>i$, $L_C$ has $u_j\to u_i$, so a disagreement means $u_i\to u_j$ in $S$, and by (i) there are finitely many such $j$; with $v_j$, $L_C$ has $u_i\to v_j$, so a disagreement means $v_j\to u_i$ in $S$, finitely many by (iii). At $v_i$: with $r$ none; with $v_j$ for $j<i$ at most $i-1$; with $v_j$ for $j>i$, $L_C$ has $v_i\to v_j$, so a disagreement means $v_j\to v_i$ in $S$, finitely many by (ii); with $u_j$ finitely many by (iii). At $r$: none. So $D(S|_C,L_C)$ is locally finite.

\emph{Step 6: assembling $L$.} Order $V(S)$ by: $x$ before $y$ if the class of $x$ is $<$ the class of $y$, or if they lie in the same class $C$ and $x$ precedes $y$ in \eqref{eq:classorder}. This is a linear order on $V(S)$; let $L$ be the corresponding transitive tournament. A vertex $x\in C$ lies in finitely many pairs of $D(S,L)$: those inside $C$ are finite in number by Step~5, and those with vertices of other classes are edges of $J(S)$ at $x$ by Step~4, of which there are finitely many because $J(S)$ is locally finite. Hence $D(S,L)$ is locally finite.
\end{proof}

\phantomsection\label{pf:Z}
\lemZ*
\begin{proof}
Let $p\in P(x,y)$, so $x\to p\to y$, and suppose for a contradiction that $Z(x,y)$ contains three distinct vertices $z_1,z_2,z_3$, so $y\to z_i\to x$. The vertices $x,y,p,z_1,z_2,z_3$ are pairwise distinct: $p\ne x,y$ and $z_i\ne x,y$ by definition, and $p\ne z_i$ because $p\to y$ while $y\to z_i$. By the pigeonhole principle two of the $z_i$, say $z_1$ and $z_2$, are oriented the same way relative to $p$.

\emph{Case $p\to z_1$ and $p\to z_2$.} Put $A=\{p,y\}$, $B=\{z_1,z_2\}$ and $c=x$. Then $p\to y$ inside $A$, so $A$ induces $\TT_2$; $B$ induces $\TT_2$ whichever way the arc between $z_1$ and $z_2$ points; $A\to B$ because $p\to z_1,z_2$ (the case hypothesis) and $y\to z_1,z_2$ (as $z_i\in Z(x,y)$); $B\to\{c\}$ because $z_i\to x$; and $\{c\}\to A$ because $x\to p$ and $x\to y$. So $\{x,y,p,z_1,z_2\}$ induces $C_3[\TT_2,\TT_2,1]=\Tstar$.

\emph{Case $z_1\to p$ and $z_2\to p$.} Put $A=\{z_1,z_2\}$, $B=\{x,p\}$ and $c=y$. Then $A$ induces $\TT_2$; $x\to p$ inside $B$; $A\to B$ because $z_i\to x$ and $z_i\to p$; $B\to\{c\}$ because $x\to y$ and $p\to y$; and $\{c\}\to A$ because $y\to z_i$. Again $\{x,y,p,z_1,z_2\}$ induces $\Tstar$.

Both cases contradict the $\Tstar$-freeness of $S$. Hence $|Z(x,y)|\le 2$.
\end{proof}

\phantomsection\label{pf:Pempty}
\lemPempty*
\begin{proof}
Let $x$ be a vertex and suppose $x\to y$ and $x\to z$ with $y\ne z$ both satisfy $P(x,y)=P(x,z)=\varnothing$. Either $y\to z$, in which case $x\to y\to z$ gives $y\in P(x,z)$, or $z\to y$, in which case $z\in P(x,y)$; both are impossible. So at most one out-arc of $x$ has empty $P$. Dually, suppose $y\to x$ and $z\to x$ with $y\ne z$ both satisfy $P(y,x)=P(z,x)=\varnothing$. Either $y\to z$, giving $y\to z\to x$ and $z\in P(y,x)$, or $z\to y$, giving $y\in P(z,x)$; both are impossible. So at most one in-arc of $x$ has empty $P$, and $x$ has degree at most $2$ in the undirected graph formed by the arcs with empty $P$.
\end{proof}

\phantomsection\label{pf:structure}
\thmstructure*
\begin{proof}
Let $xy\in E(J(S))$ with $x\to y$. Then $x\to y$ lies in infinitely many directed triangles, i.e.\ $Z(x,y)$ is infinite and in particular $|Z(x,y)|\ge 3$. By \cref{lem:Z}, $P(x,y)=\varnothing$. Hence every edge of $J(S)$ is an arc with empty $P$, and \cref{lem:Pempty} bounds the degree of every vertex in $J(S)$ by $2$. In particular $J(S)$ is locally finite, and if $S$ is countable the sufficiency direction of \cref{lem:criterion} yields a transitive tournament $L$ on $V(S)$ with $D(S,L)$ locally finite.
\end{proof}

\phantomsection\label{pf:onearc}
\lemonearc*
\begin{proof}
Let $u\to v$ be the reversed arc, so $L'$ has $v\to u$ and agrees with $L$ on every other pair. If a directed triangle of $L'$ did not use the arc $v\to u$, all three of its arcs would be arcs of $L$, and $L$ would contain a directed triangle, which is impossible. Hence every directed triangle of $L'$ contains the arc $v\to u$, and any two directed triangles of $L'$ share an arc. Now suppose $X\subseteq V(L')$ induces a copy of $\Tstar$. The triangles $c\to a_0\to b_0\to c$ and $c\to a_1\to b_1\to c$ of $\Tstar$ use the arc sets $\{ca_0,a_0b_0,b_0c\}$ and $\{ca_1,a_1b_1,b_1c\}$, which are disjoint; their images are two arc-disjoint directed triangles of $L'$, a contradiction. So $L'$ is $\Tstar$-free.
\end{proof}

\phantomsection\label{pf:main}
\thmmain*
\begin{proof}
Let $S$ be a countably infinite $\Tstar$-free tournament.

\emph{Case 1: $S$ is not transitive.} By \cref{thm:structure} there is a transitive tournament $L$ on $V(S)$ with $D(S,L)$ locally finite. Since $S$ is not transitive and $L$ is, $L\ne S$. Every subtournament of a transitive tournament is transitive, and $\Tstar$ is not, so $L$ is $\Tstar$-free. Take $S'=L$.

\emph{Case 2: $S$ is transitive.} As $S$ has infinitely many vertices it has an arc $u\to v$; let $S'$ be obtained from $S$ by reversing it. Then $S'\ne S$, $D(S,S')$ is the single edge $uv$, which is locally finite, and $S'$ is $\Tstar$-free by \cref{lem:onearc}.

In both cases $S'\ne S$ is a $\Tstar$-free tournament on $V(S)$ with $D(S,S')$ locally finite. Since $S'\ne S$, the set of reversed pairs is nonempty, so $S'$ is a locally finite perturbation of $S$ in the sense of \cite{BGJMS25}, and it is $\Tstar$-free.
\end{proof}

\section{Report of the LLM referee (verbatim)}\label{app:referee}

\begin{tcolorbox}[colback=gray!8,colframe=gray!50,boxrule=0.4pt,arc=2pt]
\small The following is the unedited report of the adversarial referee agent (\texttt{claude-opus-5}, 2026-09-17) on the \emph{original} writeup. Its step numbers (S1--S11) and lemma numbers refer to that writeup, not to this note: the writeup's Lemma~1, 2, 3 and Corollary~4 are \cref{lem:criterion,lem:Z,lem:Pempty,thm:structure} here, and its \S4 Cases~1 and~2 are the proof of \cref{thm:main} and \cref{lem:onearc}. The scripts it mentions are in \texttt{verification\_astra/scripts/2506.08810\_\_03/}. The out-degree tuple in its step~S1 has its last two entries transposed (see \cref{rem:comp}). Treat it as machine output, not as an authority.
\end{tcolorbox}

{\setlength{\emergencystretch}{3em}\input{.build/referee.tex}\par}

\begin{thebibliography}{9}
\small
\setlength{\itemsep}{2pt}
\bibitem{BGJMS25} M.~Bonamy, C.~Groenland, T.~Johnston, N.~Morrison and A.~Scott, Infinite induced-saturated graphs, \href{https://arxiv.org/abs/2506.08810}{arXiv:2506.08810} (v3, September 2025); Canadian Journal of Mathematics, FirstView (2026), \href{https://doi.org/10.4153/S0008414X26102132}{doi:10.4153/S0008414X26102132}, published online 24 March 2026. Problem~23 and Conjecture~24 are in Section~8.
\bibitem{Sz30} E.~Szpilrajn, Sur l'extension de l'ordre partiel, Fundamenta Mathematicae 16 (1930), 386--389.
\end{thebibliography}

\end{document}
